\section{Comparison with Existing Standards}
\label{sec:comparison}

\begin{table}[H]
\centering
\caption{Security token standard comparison.}
\label{tab:standard-comparison}
\begin{tabular}{lcccc}
\toprule
\textbf{Feature} & \textbf{Lux} & \textbf{ERC-3643} & \textbf{ERC-1400} & \textbf{ERC-1404} \\
\midrule
Partitioned holdings & Yes & No & Yes & No \\
Dynamic compliance oracle & Yes & Yes (ONCHAINID) & No & No \\
MPC controller & Yes & No & No & No \\
Post-quantum signatures & Yes & No & No & No \\
Forced transfers & Yes & Yes & Yes & No \\
Document management & Yes & No & Yes & No \\
DID-based identity & Yes & Yes (ONCHAINID) & No & No \\
Hierarchical partitions & Yes & No & Partial & No \\
Vesting schedules & Yes & No & No & No \\
Convertible instruments & Yes & No & No & No \\
\bottomrule
\end{tabular}
\end{table}

\subsection{ERC-3643 (T-REX)}

ERC-3643, the Token for Regulated Exchanges (T-REX) standard, is the closest competitor to the Lux standard. Both use on-chain identity registries for compliance checking. Key differences:

\begin{itemize}[nosep]
  \item T-REX uses ONCHAINID, a centralized identity provider. Lux uses I-Chain DIDs with verifiable credentials from multiple providers.
  \item T-REX has a single controller address. Lux uses MPC threshold controllers.
  \item T-REX does not support partitioned holdings or post-quantum signatures.
  \item T-REX is deployed on Ethereum mainnet with high gas costs. Lux EVM provides sub-cent transfer costs.
\end{itemize}

